\documentclass[11pt]{article}
\usepackage[margin=1in]{geometry}
\usepackage{amsmath,amssymb,hyperref}
\newcommand{\F}{\mathbb F}
\newcommand{\E}{\mathbb E}
\newcommand{\Nm}{\operatorname{Norm}}
\newcommand{\Tr}{\operatorname{Tr}}
\newcommand{\ord}{\operatorname{ord}}
\title{A single algebraic norm for the target parity and its dyadic restrictions}
\author{Project research note}
\date{18 November 2026}
\begin{document}
\maketitle

For independently hashed slots on an odd-characteristic elliptic curve,
the parity of the ordered target count is the parity of the order of
vanishing of one scalar norm. At the project's isolation widths,
the expected cost of handling the necessary Taylor precision is a
constant factor per query for a uniform division-free norm circuit.
On a draw with a unique tuple, first derivatives of the same norm
give all its abscissas without bisection. We give a finite exact
evaluator and price its dense algebra and an alternative balanced
elimination. Neither gives a new end-to-end exponent.

\section{Eliminate the signs without losing their multiplicities}
Write the curve as $y^2=x^3+ax+c$ with $p>3$. Fix a target with
$2T\ne O$, and put $t=x(T)$. First suppose all slot points come
in pairs of distinct opposites. In slot $i$ let $U_i$ contain the
accepted abscissas that actually lift to two rational points, and let
$g_i(X)=\prod_{u\in U_i}(X-u)$. These sets can be built at the
usual individual slot cost. An empty slot gives count zero. Use
\[
  A=\F_p[X_1,\ldots,X_m]/(g_1(X_1),\ldots,g_m(X_m)),\qquad
  D=\dim A=\prod_i |U_i|.
\]
The Chinese remainder theorem identifies $A$ with one copy of $\F_p$
per abscissa tuple. Monomial coordinates permit joint arithmetic before
choosing any such tuple.

Let $K(u,v,z)$ denote the third summation polynomial
\[
 (u-v)^2z^2-2\big((u+v)(uv+a)+2c\big)z
             +(uv-a)^2-4c(u+v).
\]
Its roots on the projective line are $x(P+Q),x(P-Q)$ when
$u=x(P),v=x(Q)$, with $x(O)=\infty$. For different finite
abscissas this follows from the two chord slopes. The bihomogeneous
continuation in the inputs gives $(Z_0-vZ_1)^2$ if one input is
infinity and the other is $v$. For equal finite inputs $u=v$ the
leading coefficient vanishes and the middle coefficient is
$-4(u^3+au+c)$. This gives one infinity and the usual doubling
root when $P$ is nonfixed. At rational two-torsion the constant
coefficient is $(3u^2+a)^2\ne0$, giving two infinities.
Both inputs at infinity also give two infinities. This is the
summation-polynomial recurrence of
\href{https://eprint.iacr.org/2004/031}{I. Semaev}, used here with
projective degrees and multiplicities.

Start with $S_1(Z)=Z-X_1$. For $j>1$ take the binary homogeneous
resultant in $W$ of $S_{j-1}(W)$ and $K(W,X_j,Z)$, preserving the
formal degrees $2^{j-2}$ and $2$ in $W$. Denote the result by
$S_j(Z)$, of formal degree $2^{j-1}$ in $Z$. For every grid
specialization, this form has a root, with multiplicity, for each
choice of the first $j$ signs modulo simultaneous reversal.
Indeed, a representative of each previous sign orbit extends in
two ways and the two new abscissas are the roots of its $K$ factor.
If the old sum is fixed under negation, the two new values coincide
and their multiplicities add. The homogeneous resultant multiplies
these factors, including those from infinite roots. Every specialized
form is nonzero because each individual quadratic is nonzero.

Set $S=S_m$ and let $S_{\mathbf u}$ be its specialization at a grid
point. Each sign orbit with final abscissa $t$ contains exactly one
tuple summing to the specified signed target $T$, since $2T\ne O$.
The resulting integer identity also covers $H(T)\ge p$.
\begin{equation}
 H(T)=\ord_{\epsilon=0}\mathcal N(\epsilon),\qquad
 \mathcal N(\epsilon)=\Nm_{A/\F_p}\big(S(t+\epsilon)\big)
     =\prod_{\mathbf u}S_{\mathbf u}(t+\epsilon). \label{eq:norm}
\end{equation}
The norm is the determinant of multiplication on $A$. The product
proves the identity. It need not be the chosen evaluation algorithm.
For a subwindow replace the $g_i$ by the corresponding divisors.

If slots contain rational two-torsion, expand each slot into its
ordinary paired part and up to three individual fixed points. In
each branch subtract the sum of the fixed choices from $T$. The
result still satisfies $2T'\ne O$. Apply (\ref{eq:norm}) at
$x(T')+\epsilon$ to its ordinary slots. The all-fixed branch
contributes no tuple. Multiplying the nonempty branch norms adds
their orders. There are at most $4^m$ branches, a lower-order
logarithmic cost in the project regime.

\section{Cost of precision and recovery of an isolated tuple}
Use the tuned widths of \texttt{../mixed\_coefficient/mixed.tex}.
Let $H_0$ count tuples at the original window. That note proves
$\E H_0=O(1)$, $\E[H_0(H_0-1)]=O(1)$ and a probability at least
$2/9-o(1)$ of exactly one tuple, passing the chord guards.
Every nested restriction satisfies $H'\le H_0$ deterministically.
Evaluate its norm modulo $\epsilon^k$, doubling $k$ from one until
a nonzero coefficient is found. Its least index is the exact
integer $H'$. In particular no large-characteristic lifting
assumption is needed to read its parity.

Suppose a division-free circuit of size at most $L$ computes the
scalar norm with its input coefficients, uniformly over these
widths, and is valid over truncated polynomial rings. A ring
operation at precision $k$ costs $O(k^2)$ base-field operations by
ordinary multiplication. The sum over doublings is $O((1+H_0)^2)$.
Since $\E(1+H_0)^2=1+3\E H_0+\E[H_0(H_0-1)]=O(1)$, the expected
cost of an adaptive sequence of at most $Q$ queries is $O(QL)$,
apart from obtaining slot inputs and constructing circuits. This
argument uses a uniform cost bound $L$ and the original seed law.
It makes no independence assumption on the adaptive child queries.
It does not give a bound of the desired size for $L$.

There is a cheaper interface when only a relation is needed.
Introduce scalar perturbations $s_0,\ldots,s_m$ in the norm
\[
  N(z,\mathbf s)=\Nm\big(S(z)+s_0+\textstyle\sum_i s_iX_i\big)
\]
and evaluate its value and first derivatives at $(t,\mathbf0)$.
A nonzero value means $H=0$. A zero value with $\partial_z N\ne0$
means $H=1$. In the latter case only one grid point $\mathbf u^*$
is a zero. Differentiating the split product gives
\[
  \partial_{s_0}N=\prod_{\mathbf u\ne\mathbf u^*}S_{\mathbf u}(t)\ne0,
       \qquad u_i^*=\frac{\partial_{s_i}N}{\partial_{s_0}N}.
\]
Check at most $2^m$ signs and the original guards after this readout.
At $H\ge2$ the value and $z$ derivative both vanish and the relation
finder may abstain. Thus root derivatives suffice at the previous
isolation probability without adaptive queries. With fixed-point
branches, use their product and perturb each factor with its own
corresponding abscissas, including constants in the fixed slots.
First-order forward differentiation costs $O(m)$ field operations
per circuit gate. The finite verifier uses the cofactor formula
for a determinant derivative, which works also at a singular matrix.

\section{One projector answers every restriction}
There is a second exact implementation which keeps the root algebra.
Write $S(t+\epsilon)=\sum h_j\epsilon^j$ and $d=2^{m-1}$. Define
\[
  e_j=\prod_{k=0}^{j-1}(1-h_k^{p-1})\quad(1\le j\le d),\qquad
  B=\sum_{j=1}^d(-1)^{j-1}e_j .
\]
On each grid point $e_j$ indicates root order at least $j$, so $B$
is the zero-one indicator of odd local multiplicity. Since every
specialization is a nonzero form of degree $d$, this includes every
multiplicity. Equivalently, for multiplication operators $M_{h_k}$
the integer count, before any reduction, is
\[
 H(T)=\sum_{j=1}^{d}\dim_{\F_p}\bigcap_{k<j,\ k\ge0}\ker M_{h_k}.
\]
Indeed the $j$th subspace is precisely the image of $e_j$. The first
kernel has dimension at most $H(T)$, and dimension one on an isolated
draw. Extracting this small kernel from a succinct description of
the multiplication operator is itself a computational problem.
For subwindow sets $V_i\subseteq U_i$ let $I_i$ be
the univariate Lagrange indicator of $V_i$ on $U_i$. Then
\begin{equation}
   (-1)^{H_V(T)}
   =\Nm_{A/\F_p}\left(1-2B\prod_i I_i(X_i)\right). \label{eq:bit}
\end{equation}
The right side is computed over the odd prime field and is exactly
$+1$ or $-1$. The formula holds for arbitrary seed choices and
simultaneously for all subwindows. On an isolated root $B$ consists
of one tensor product of univariate Lagrange indicators. In this
case $\Tr(X_iB)=u_i^*$.

On independent root seeds, replacing $B$ by just $e_1$ gives the
same answer for every restriction except on an event of probability
at most $\lambda |G[2]|(2^m-2)/|G|$. Here $\lambda=\prod_i
(|G|\tau_i/2^b)/|G|$. Indeed, a difference requires two distinct
signed tuples with the same abscissas reaching $T$ at the original
window. The all-abscissas-equal term in the previous factorial
moment bound counts precisely such ordered pairs. Restricting
windows cannot create a pair. Equation (\ref{eq:bit}) handles the
exceptional seeds exactly.

\section{Pricing the joint operation}
Consider prime-order curves of order $r$ and tuned balanced widths.
Put $\rho_i=\tau_i/2^b$, $\nu_i=r\rho_i$ and $n=(r-1)/2$.
Independence and the nonpositive covariance for distinct abscissas
in each slot give
\[
 \E D=\prod_i n\rho_i=\lambda r(1-1/r)^m/2^m,\qquad
 \frac{\operatorname{Var} D}{(\E D)^2}
     \le\prod_i(1+1/(n\rho_i))-1=o(1).
\]
The last estimate uses $m/\min_i\nu_i=o(1)$. Hence a dense root
element costs $2^{b-m+O(1)}$ field entries with probability $1-o(1)$.
The companion-matrix verifier spends additional space on matrices.
Even using dense field elements alone costs $2^{(1-o(1))b}$ for
the growing polylogarithmic arities here. The constant-probability
isolation event intersects this typical size with the same
asymptotic conditional probability. Small final projector rank
alone gives no algorithm for obtaining its factors.

A balanced elimination could instead build squarefree univariate
polynomials for the two halves, dropping every even-multiplicity
root. The following estimates price this explicit representation.
For a fixed block of $s$ slots, put $A_s=\prod_{i\text{ in block}}\nu_i$
and let $h(U)$ be its ordinary count at $U$. Summing over $U$ gives
\[
 \E\sum_U h(U)(h(U)-1)
 \le \frac{A_s^2}{r}\prod_{i\text{ in block}}(1+2/\nu_i)
       +\frac{A_s|G[2]|(2^s-1)}r .
\]
For a proper set of matching abscissas choose the first tuple in
at most $r^s$ ways, allow two signs at matched positions and force
the last unmatched point by equal sums. For all abscissas matched,
a nonempty flip set must sum into $G[2]$, giving at most
$|G[2]|r^{s-1}$ first tuples per flip set. This includes flipping
the whole block, since the block target varies. Hash retention
probabilities then give the displayed terms.
The total number of tuples concentrates at $(1+o(1))A_s$ in relative
probability by the same product-variance calculation as for $D$.
The loss on reducing all counts modulo two is at most
$\sum_U h(U)(h(U)-1)$. For a fixed balanced half with $A_s=o(r)$
and $2^s=o(r)$ the expected loss is $o(A_s)$. Thus the number
of distinct odd abscissas is $(1+o(1))A_s/2$ in probability.
For $s=\lfloor m/2\rfloor$ and growing $m$ this is
$2^{(1/2+o(1))b}$. It prices expanding the half polynomial into
a dense array or root list. An implicit scalar norm could have a
different cost.

The existing conditional end-to-end estimate consequently remains
\[
 \log_2\E T_{\rm solve}\le
 (26.2933775156+o(1))\,b/(\log_2 b)^{1/4}
\]
under the finite-packet assumption and inherited hypotheses in
\texttt{../short\_hash\_count/count.tex}. A smaller uniform circuit
for the scalar norm, in particular for its value and first derivatives
at the target, would address the remaining joint computation.

\paragraph{Finite verification.}
Run \texttt{python -B research/joint\_parity/verify.py} with NumPy.
The checker constructs the slot companion matrices in a monomial
basis. For a binary form $P(W)=\sum p_iW^i$ of formal degree $h$,
its resultant with $AW^2+BW+C$ can be computed without divisions as
\[
 \sum_{i=0}^{h}p_i A^{h-i}
   \left(p_i C^i+\sum_{j<i}p_j C^j L_{i-j}\right),\qquad
 L_0=2,\ L_1=-B,\ L_k=-B L_{k-1}-AC L_{k-2}.
\]
This follows by expanding $A^hP(w_1)P(w_2)$ at the two formal roots
and continues polynomially to $A=0$. The checker interpolates the
output in $Z$ on fixtures with $p>d$. It checks the exact integer
orders via kernel dimensions and the norm signs against independent
elliptic addition. Reuse of the root projectors passes
567 child/target tests over all 63 dyadic boxes on the curve of order
19. Both trace readout and differentiated norm recover all 76
isolated cases in those boxes. Four repeated slots give counts as
large as 28 in characteristic 13. Fixed-point branching passes all
four eligible signed targets on the order-8 curve with full rational
two-torsion.

On the order-65419 fixture over $\F_{65521}$, four independent
seed draws (one per slot) at widths $(8,16,16,16)$ gave a grid of
dimensions $(4,7,10,10)$ at a verified singleton target $(0,1)$.
Across the $28$-by-$100$ split the field rank of $S(t)$ was 15,
the rank of its square was 28, and every sampled successive square
through exponent $2^{15}$ also had rank 28. The final zero projector
had rank one. These are exact field ranks for one disclosed fixture.
They test the proposal to keep ranks small during powering and do
not supply an asymptotic lower bound for other circuits.
\end{document}
